\subsection{The construction is the only shape}\label{sec:characterisation}

Section~\ref{sec:extremal} found the shape of Theorem~\ref{thm:construction} in every enumerated extremal graph with
$s(G)=1$. We now show that, given Didin's main theorem, no other shape exists at any order. The argument has two
parts: equality in the proof of Proposition~\ref{prop:where}(a), together with Didin's inequality applied to the
block, reduces the block to a graph on at most seven vertices with prescribed size; and an exhaustive check over all
graphs on five to seven vertices leaves exactly $X_6$ and $X_7$, at exactly the attaching triples observed.

\begin{lemma}[conditional on {\cite[Theorem~3]{Didin}}]\label{lem:block-reduction}
Let $(n,\alpha,3)$ be in the second regime with $\alpha\ge3$, and let $G$ be $3$-connected with $n$ vertices,
independence number $\alpha$, $t(n,\alpha)+\ceil{3\alpha/2}$ edges, $p(G)\ge2$ and no partition into $\alpha$
cliques. Let $H$ be an $\alpha$-core of $G$ with the largest number of components. Then:
\begin{enumerate}
\item[(i)] $H$ has $\alpha-1$ components; exactly one of them, $X$, has $\alpha(X)=2$, the others are cliques; every
component induces in $G$ the same graph as in $H$; and every component sends exactly three edges to the others;
\item[(ii)] the clique components have orders $3$ and $4$, and $X$ has $m\in\{5,6,7\}$ vertices and exactly
$t(m+1,2)-1\in\{5,8,11\}$ edges;
\item[(iii)] every component has exactly three vertices with a neighbour in another component, each with exactly one
such neighbour, so the edges between components form a matching, and the quotient $Q$ on the components is a loopless
cubic multigraph that is $3$-edge-connected;
\item[(iv)] with $A$ the three attaching vertices of $X$: every vertex of $X-A$ has degree at least $3$ in $X$, every
vertex of $A$ degree at least $2$, every component of $X-T$ meets $A\setminus T$ for every set $T$ of at most two
vertices of $X$, and $\theta(X)\ge3$.
\end{enumerate}
\end{lemma}

\begin{proof}
(i) By Proposition~\ref{prop:where}, $\alpha$ is odd, $s(G)=1$ and $\floor{n/\alpha}=3$, so $H$ has $p=\alpha-1$
components $V_1,\dots,V_{\alpha-1}$. Let $q$ be the number of edges of $G$ between components. The proof of
Proposition~\ref{prop:where}(a) gives the chain
\[
\e(G)\ \ge\ \e(H)+q\ \ge\ t(n+1,\alpha)-1+\tfrac{3\alpha-3}{2}\ =\ t(n,\alpha)+2+\tfrac{3\alpha-3}{2}\ =\ \e(G),
\]
using \eqref{eq:didin} for $H$, $|\partial V_i|\ge\lambda(G)\ge\kappa(G)=3$ with $\sum_i|\partial V_i|=2q$, and
Lemma~\ref{lem:tshift}. Every inequality is therefore an equality: $\e(H)=t(n+1,\alpha)-1$, every component sends
exactly three edges, and $G$ has no edge inside a component outside $H$. The independence numbers of the
$\alpha-1$ components sum to $\alpha(H)=\alpha$ and are at least one each, so exactly one is two; a connected graph
with independence number one is complete.

(ii) $X$ is $\alpha$-critical, because deleting an edge of $X$ that keeps $\alpha(X)$ would keep $\alpha(H)$; so the
only spanning subgraph of $X$ with independence number two is $X$ itself, and $p(X)=1$, $s(X)=1$. Didin's Theorem~3
for $X$, whose hypothesis $m\ge\max(\alpha(X)+1,p(X))=3$ holds, gives $\e(X)\ge t(m+1,2)-1$. With
$a_1,\dots,a_{\alpha-2}$ the clique orders,
\[
t(n+1,\alpha)-1\ =\ \e(H)\ \ge\ \sum_i\binom{a_i}{2}+t(m+1,2)-1\ \ge\ t(n+1,\alpha)-1,
\]
because the sum counts the edges of $\alpha$ cliques on $n+1$ vertices. Equality forces $\e(X)=t(m+1,2)-1$ and all
$\alpha$ clique orders $a_1,\dots,a_{\alpha-2},\floor{(m+1)/2},\ceil{(m+1)/2}$ to be balanced over $n+1$ vertices,
hence in $\{3,4\}$ since $3\alpha+2\le n+1\le4\alpha$; so $m+1\in\{6,7,8\}$.

(iii) Every component has at least three vertices by (ii). Let $A_P$ be the vertices of a component $P$ with a
neighbour outside $P$. If $|A_P|\le2$, then $A_P$ separates $P-A_P\neq\emptyset$ from the other components, a vertex
cut of size at most two. So $|A_P|=3$, and by (i) each vertex of $A_P$ carries exactly one of the three edges. A cut of
at most two edges of $Q$ yields, by taking one end of each cut edge, a vertex cut of $G$ of size at most two with both
sides nonempty, so $Q$ is $3$-edge-connected.

(iv) A vertex of $X-A$ has all its neighbours in $X$, and $\delta(G)\ge3$; a vertex of $A$ has exactly one neighbour
outside. If a component of $X-T$ misses $A\setminus T$, then $T$ separates it from the other components. If
$\theta(X)\le2$, the clique components and two cliques covering $X$ partition $V(G)$ into $\alpha$ cliques.
\end{proof}

The conditions of Lemma~\ref{lem:block-reduction} on $(X,A)$ are finite. We tested them on every graph with five, six
or seven vertices and every set $A$ of three of its vertices (EXP-014, \texttt{check\_blocks.py}, NetworkX's graph
atlas, whose \nblockatlas\ graphs on these orders agree with the counts of OEIS A000088 \cite{OEIS}). Independence
and clique cover numbers were computed by two methods each, with no disagreement. Table~\ref{tab:blockcheck} gives
the number of survivors after each condition.

\begin{table}[t]
\centering\small
\caption{The finite check behind Proposition~\ref{prop:blockcheck}: graphs on five to seven vertices and attaching
triples that survive each necessary condition of Lemma~\ref{lem:block-reduction}, and the isomorphism classes of the
marked blocks that survive all of them, with the observed triple each class equals.}
\label{tab:blockcheck}
\begin{tabular}{lr}
\toprule
condition & survivors \\
\midrule
graphs on $5$ to $7$ vertices & 1234 \\
C1: order and size & 178 \\
C2: connected, $\alpha=2$, every edge critical & 3 \\
C3: $\theta\ge3$ & 3 \\
C4: degrees, pairs $(X,A)$ & 39 \\
C5: components of $X-T$, pairs $(X,A)$ & 18 \\
\midrule
surviving class (order, members) & equals \\
\midrule
$m=6$, 2 triples & $X_6$ $\{0, 1, 5\}$ \\
$m=7$, 4 triples & $X_7$ $\{1, 2, 6\}$ \\
$m=7$, 4 triples & $X_7$ $\{1, 2, 3\}$ \\
$m=7$, 8 triples & $X_7$ $\{0, 1, 2\}$ \\
\bottomrule
\end{tabular}

\end{table}

\begin{proposition}[machine-verified]\label{prop:blockcheck}
A pair $(X,A)$, with $X$ a graph on $m\in\{5,6,7\}$ vertices and $t(m+1,2)-1$ edges that is connected, has independence number $2$ and only critical edges, and $A$ a set of three of its vertices, satisfies the conditions of Lemma~\ref{lem:block-reduction}(iv) if and only if $X\cong X_6$
with $A$ in the automorphism orbit of $\{0,1,5\}$, or $X\cong X_7$ with $A$ in one of the orbits of $\{0,1,2\}$,
$\{1,2,3\}$ and $\{1,2,6\}$, in the labelling of Figure~\ref{fig:blocks}.
\end{proposition}

\begin{theorem}[conditional on {\cite[Theorem~3]{Didin}}]\label{thm:characterise}
Let $(n,\alpha,3)$ be in the second regime with $\alpha\ge3$. A $3$-connected graph with $n$ vertices, independence
number $\alpha$, $t(n,\alpha)+\ceil{3\alpha/2}$ edges and $p(G)\ge2$ has no partition into $\alpha$ cliques if and only
if it is a graph of Theorem~\ref{thm:construction}: $\alpha$ is odd, $3\alpha<n\le4\alpha-1$, $V(G)$ splits into the
cliques of $T(n,\alpha)$ with two cliques replaced by $X_6$ (two of order $3$) or $X_7$ (one of order $3$ and one of
order $4$), $X$ attaches at a triple of Proposition~\ref{prop:blockcheck}, and the edges between parts form a matching
along a $3$-edge-connected loopless cubic quotient.
\end{theorem}

\begin{proof}
The graphs of Theorem~\ref{thm:construction} have the stated properties, unconditionally. Conversely, by
Lemma~\ref{lem:block-reduction} and Proposition~\ref{prop:blockcheck}, $X$ is $X_6$ or $X_7$ at an observed triple and
the clique orders lie in $\{3,4\}$ and sum to $n-|X|$. With $n=3\alpha+r$, $X_6$ leaves $r$ cliques of order $4$ and
$\alpha-2-r$ of order $3$, which are the cliques of $T(n,\alpha)$ without two of order $3$; $X_7$ leaves $r-1$ and
$\alpha-1-r$, the cliques of $T(n,\alpha)$ without one of order $3$ and one of order $4$. Part (ii) of the lemma gives
the matching and the quotient, and Proposition~\ref{prop:where} gives $n\le4\alpha-1$.
\end{proof}

\begin{corollary}[conditional on {\cite[Theorem~3]{Didin}}]\label{cor:characterise}
The extremal graphs of the second regime without a partition into $\alpha$ cliques are exactly the graphs of
Theorem~\ref{thm:construction}, at the triples where they are extremal ($n\le24$ by Theorem~\ref{thm:values}, and every
admissible $n$ if Didin's Theorem~9 holds), and the \ntenthreeb\ graphs at $(10,3,3)$.
\end{corollary}

\begin{proof}
By Proposition~\ref{prop:where}, such a graph is in case (a), covered by Theorem~\ref{thm:characterise}, or in case
(b) at one of six triples, all with $n\le16$, where the enumeration of Section~\ref{sec:extremal} lists every extremal
graph and finds case (b) only at $(10,3,3)$.
\end{proof}

All four structural claims that the data suggested rest on Didin's Theorem~3: the matching and the cubic, $3$-edge-connected quotient through the equalities of Proposition~\ref{prop:where}(a), and the clique orders and the two blocks through his theorem applied once more, to the block.
